\chapter{A bounded-shape one-weak C2 family, at all tests}
\label{component:pec59340d644f}
\noindent\textbf{Component }\path{pec59340d644f}\par
{\small\noindent Source: \path{R13/baseline/baseline/baseline/baseline/baseline/baseline/advances/c2/general_single_angle.tex}\par}
{\small\noindent SHA256: \path{ec59340d644fc7d091e3d6cf8f9da2f1a1096c7fb209e9c0a8233c82509a782f}\par}
{\small\noindent Inventory use: Complete six-term Laurent identity, full scalar and both normal monic divisions, literal n=0 corner and all four outer-core identities\par}
\medskip
\section{A bounded-shape one-weak C2 family, at all tests}
\label{pec59340d644f:sec:c2-general-single-angle}

This module gives a source-domain extension of the fixed rational
family.  Its scalar uniqueness proof is an exact six-term Laurent
identity and a three-cycle arithmetic--geometric mean inequality.
The full-normal and outer-chart arguments include literal $n=0$.

Put $x>0$ infinitesimal in the coefficient field and $\rho=\sqrt{x}$.
Choose positive bounded coefficient units $k,b,r,h$ such that
\[
 1-r,\quad b-h,\quad g=b-1+r+xh,\quad b-k^2r
\]
are positive bounded units.  Set
\[
 c=\rho,\quad a=k/\rho,\quad t=x(b-h)/k^2,
 \quad j=1+xr-t,\quad a_2=k^2/x,
 \quad A_c=1+xr,\quad S_c=A_c+a_2xt.
\]
Thus $h=b-a_2t$, $g=(1+x)b-(1-r+a_2xt)$, and $j=1+xr-t$.
Use the exact original rational source
\begin{align*}
 \Delta={}&(1+x)xb^2+2x(1+x)br+(1+x)(1+a_2x)b
       +x^2(1-a_2)r^2+2x(1+a_2x)r+1+a_2x,\\
 N={}&(1+x)(b+r)(b+a_2(1-t))-a_2(1+a_2x)(1-r)t,\\
 n={}&\frac{(b+r)A_cS_c-r\Delta}{2r\Delta},\\
 m={}&\frac{(1+a_2x)(1-r)A_cS_c}{2g\Delta},\\
 C={}&\frac{A_cN}{2a_2t\Delta},\\
 A={}&\frac{a_2A_c(b+r)S_c}{2h\Delta},\\
 D={}&\frac{(1+x)(b+r)(b-a_2xr)S_c}{2a_2j\Delta},\\
 \alpha={}&\frac{x^2(b+r)N}{2(1-xh)\Delta}.
\end{align*}
Assume $n\geq0$ and physicality of this canonical source.  The latter
is the same inherited source-physicality hypothesis used for the
general noncommuting-residue exit; no conclusion below is predicated
on a new source enumeration.  Literal $n=0$ is allowed.

The parent, with $d=1$ and $\lambda=\beta=0$, is
\begin{align*}
 F={}&4nD_uD_vX^2+4mM^2+4CD_v(uW-Y)^2+4AD_uD_vY^2
                +4DJ^2+4\alpha D_u(vY-Z)^2+H,\\
 M={}&vW-X+uvY-uZ,\qquad J=uW+uvX-Y-vZ,\\
 H/2={}&J^2+D_u(vY-Z)^2+D_v(X+uZ)^2-D_uD_vZ^2,
 \qquad D_u=1+u^2,\quad D_v=1+v^2.
\end{align*}
Its exact stationary contact is
\[
 (u,v;W,X,Y,Z)=(\rho,k;\rho b/k,-\rho r,xh/k,1).
\]
Writing $\kappa=b/k$ and using a subscript zero for residues gives
\begin{align*}
 \Delta_0&=(1+k_0^2)(b_0+1),\qquad
 n_0=\frac{b_0+r_0-r_0\Delta_0}{2r_0\Delta_0}\geq0,\\
 m_0&=\frac{1-r_0}{2(b_0-1+r_0)(b_0+1)}>0,\\
 C_*:=(xC)_0&=\frac{k_0^2(b_0+r_0)}{2(b_0-h_0)\Delta_0},\qquad
 A_*:=(xA)_0=\frac{k_0^2(b_0+r_0)}{2h_0\Delta_0},\\
 D_*:=(D/x)_0&=\frac{(b_0+r_0)(b_0-k_0^2r_0)}{2k_0^2\Delta_0},\qquad
 E_*:=(\alpha/x)_0=\frac{k_0^2(b_0+r_0)^2}{2\Delta_0}.
\end{align*}
All starred quantities are positive.  Thus $n$ is bounded, $m$ is a
positive unit, $C,A$ have scale $x^{-1}$, and $D,\alpha$ have scale $x$.

\begin{theorem}[General bounded-shape single-angle source]
\label{pec59340d644f:thm:c2-general-single-angle}
Every source in the displayed domain has tensor SOS for the full
original parent at all independent real or complex projective testing
states and all projective testing angles.  No comparison between a
coefficient infinitesimal and a testing infinitesimal is needed.
The conclusion includes $n=0$.  Applying the exact C2 reflection gives
the corresponding original orientation with weak $C$ and strong
$\alpha$; this is a sufficient source domain, not a global one-weak
C2 classification.
\end{theorem}

\paragraph{The finite $u=0$ axis and its full normals.}
Set $a_0=4m_0$, $b_0'=a_0+2$, and
\[
 H_0(v)=(4n_0+2)(1+v^2)+a_0,\quad \eta=h/b,
 \quad \gamma=-r-\frac{a_0b-b_0'}{H_0(k)}.
\]
Here $\gamma$ is infinitesimal and $\eta_0=C_* /(C_*+A_*)$.
Use the exact contact graph
\[
 X_g=\frac{a_0vW-b_0'uZ}{H_0(v)}+\gamma uZ,
 \qquad Y_g=\eta uW.
\]
The denominator $H_0(v)$ is a pure testing unit and $H_0(k)$ is a
pure coefficient unit.  Let $N_X=X-X_g$, $N_Y=(Y-Y_g)/\rho$.
Writing $f=F|_g=f_{ww}W^2+f_{wz}WZ+f_{zz}Z^2$, exact monic divisions give
\[
 f=aP^2+b_s B^2+c_s V^2+eE^2+g_sBV,
 \qquad(P,B,V,E)=(uW/\rho,uZ,vW,\rho Z),
\]
where
\[
 \begin{gathered}
 a=x(f_{ww}-f_{ww}|_{u=0})/u^2,\qquad
 b_s=(f_{zz}-f_{zz}|_{u=0})/u^2,\\
 c_s=f_{ww}|_{u=0}/v^2,\qquad e=f_{zz}|_{u=0}/x,\qquad
 g_s=f_{wz}/(uv).
 \end{gathered}
\]
These displayed quotients are polynomial divisions, not inverses of
testing infinitesimals.  Parity gives the divisions by $u^2,v^2,uv$;
$f_{zz}|_{u=0}=4Dv^2+4\alpha$ gives the $x$ factor.  Only $C,A$ can
produce coefficient poles, and they occur in $f_{ww}$ with a factor
$u^2$; its prefactor $x$ removes them.

Put $Q=4C_*A_* /(C_*+A_*)>0$.  The bounded coefficient residues are
\begin{align*}
 q(v)=a_0^{\rm sc}&=Q D_v,&
 b(v)&=b_0'\frac{H_0-a_0-2}{H_0},\\
 \bar c(v)&=a_0\frac{H_0-a_0}{H_0},&
 g(v)&=-2b_0'\frac{H_0-a_0}{H_0},&
 e(v)&=4D_*v^2+4E_*.
\end{align*}
Thus $c(v):=v^2\bar c(v)$ and $d(v):=-vg(v)$ give the weighted
leading scalar
\[
 q(v)u^2W^2/x+b(v)u^2Z^2+c(v)W^2-d(v)uWZ+e(v)xZ^2.
\]
The full expansion is
\[
 F=f+N_XL_X+N_YL_Y+H_{XX}N_X^2+2H_{XY}N_XN_Y+H_{YY}N_Y^2.
\]
Its hard residue is $\operatorname{diag}(H_0(v),4D_v(C_*+A_*))$.
Direct differentiation gives
\[
 L_X=V\ell_w+B\ell_z,\qquad
 L_Y=P(B_2+xe_w)+B(\rho u e_z)+E e_*,
\]
where
\[
 B_2=8D_v\{(xA)\eta D_u-(xC)(1-\eta)\},\qquad
 e_*=8v(D-\alpha).
\]
All $\ell_w,\ell_z,e_w,e_z$ are bounded; the residues of
$\ell_w,\ell_z,B_2,e_*$ vanish identically in finite $v$.
For the second identity, after subtracting $uWB_2/x$ from $F_Y|_g$,
the remaining $W$ coefficient is $ue_w$ and the remaining $Z$
coefficient is $u^2e_z+e_*$.  The constant signed-core derivatives
cancel.  This proves full normal restoration wherever the scalar
Gram on these four roots is strict.

\paragraph{Global scalar strictness by an exact cycle identity.}
Let $\nu=2n_0+1\geq1$, $\mu=2m_0+1>1$, $z=v^2$,
$\mathcal H=\nu(1+z)+\mu-1$, $e_2=4D_*$ and $e_0=4E_*$.  Then
\[
 b=\frac{2\mu(\nu(1+z)-1)}{\mathcal H},\quad
 c=\frac{2(\mu-1)z\nu(1+z)}{\mathcal H},\quad
 d=\frac{4\mu v\nu(1+z)}{\mathcal H},\quad
 e=e_2z+e_0.
\]
For $s,z>0$ define
$\mathcal D(s,z)=(q s+c)(b s+e)-d^2s/4$.
The following exact identity is the key:
\begin{align*}
 \frac{\mathcal H\mathcal D(s,z)}{(1+z)sz}
 ={}&2Q\mu(\nu-1)s/z+2Q\mu\nu s
 +2(\mu-1)\nu e_2z/s+2(\mu-1)\nu e_0/s\\
 &+Qe_0(\nu+\mu-1)/z+Qe_2\nu z\\
 &+Q\{e_2(\nu+\mu-1)+e_0\nu\}-4\mu\nu.
\end{align*}
Every nonconstant coefficient is positive, except the first may be
zero when $n_0=0$.  Exact contact and stationarity, followed by taking
residues, give a stationary zero at $(s,z)=(1,k_0^2)$.  Here is an
algebraic proof that it is the unique zero on the whole positive
quadrant.  Rescale $z=k_0^2t$ and write the Laurent expression as
\[
 \mathcal L=a_\ell s/t+b_\ell s+c_\ell t/s+d_\ell/s
                      +e_\ell/t+f_\ell t-\Sigma,
 \qquad \Sigma=a_\ell+b_\ell+c_\ell+d_\ell+e_\ell+f_\ell.
\]
The constant term has this form by its zero value at $(1,1)$.
Stationarity gives
$a_\ell-c_\ell=d_\ell-b_\ell=f_\ell-e_\ell=:\tau$.
Put $p_\ell=\min(a_\ell,c_\ell)\geq0$,
$q_\ell=\min(b_\ell,d_\ell)>0$,
$r_\ell=\min(e_\ell,f_\ell)>0$.  Exactly,
\begin{align*}
 \mathcal L={}&p_\ell(s/t+t/s-2)+q_\ell(s+1/s-2)
                              +r_\ell(t+1/t-2)\\
 &+\begin{cases}
 \tau(s/t+1/s+t-3),&\tau\geq0,\\
 (-\tau)(t/s+s+1/t-3),&\tau\leq0.
 \end{cases}
\end{align*}
Each bracket is nonnegative by the two- or three-term
arithmetic--geometric mean inequality.  The two strictly positive
pair terms force $s=t=1$ at equality.  They also give a positive
definite Hessian there: the other brackets are nonnegative and
vanish at that point.  Thus both remote uniqueness and local
nondegeneracy are proved by the same finite algebraic identity, over
any real closed residue field.  No inference from a merely local
Hessian is being used.

Consequently, for finite nonzero $v$ with $v^2\ne k_0^2$,
\[
 |d(v)|<2\sqrt{b(v)c(v)}+2\sqrt{q(v)e(v)}.
\]
Indeed the minimum of $\mathcal D(s,z)/s$ over $s>0$ is
$(\sqrt{bc}+\sqrt{qe})^2-d^2/4$ and is attained, so positivity is
strict.  Split the exact cross term using $BV=vPE$ and the pure
testing function
\[
 \theta=\frac{|v|\sqrt{b(v)\bar c(v)}}
 {|v|\sqrt{b(v)\bar c(v)}+\sqrt{q(v)e(v)}}.
\]
The resulting residue blocks on $(B,V)$ and $(P,E)$ are strictly
positive.  The denominator is a pure testing unit.  At $v_0=0$ the
same conclusion holds when $n_0>0$, with zero cross residues.
The case $n_0=0,v_0=0$ is treated separately below.

\paragraph{The matched finite-angle residue.}
For $v_0=k_0$, normalize $Z=1$ when $Z\ne0$.  If $W$ is
infinitesimal, put
\[
 R=(uW-\kappa x)/\rho,\quad S=u-\kappa^{-1}W,
 \quad U=\rho(v-k),\quad T=u(v-k).
\]
The exact conductors are
\[
 u^2-x=uS+\kappa^{-1}\rho R,\qquad
 W^2-\kappa^2x=-\kappa WS+\kappa\rho R,
 \qquad uW-\kappa x=\rho R.
\]
Write $\xi=u^2,\omega=W^2,p=uW$ and
\[
 f=p^2Q_b/x+B_b,\qquad
 Q_b=4D_v\{(xC)(1-\eta)^2+(xA)\eta^2D_u\}.
\]
Here $B_b$ is bounded and polynomial in these three invariants with
only the separate unit denominators already specified.  Put
$G=\kappa^2xQ_b+B_b$ and take successive monic differences
\[
 G=G_0(v)+(\xi-x)G_\xi+(\omega-\kappa^2x)G_\omega
                    +(p-\kappa x)G_p.
\]
Exactly,
\[
 f=R^2Q_b+\kappa^2S^2G_\omega+\rho R\mathcal A+uS\mathcal B+G_0,
\]
where $\mathcal A=2\kappa Q_b+\kappa^{-1}G_\xi+
\kappa G_\omega+G_p$ and $\mathcal B=G_\xi-\kappa^2G_\omega$.
The contact Jacobian of $(R,S)$ with respect to $(u,W)$ has
determinant $-2$.  Exact stationarity therefore gives
$\mathcal A_c=\mathcal B_c=0$.  Cancelling $\rho$ here proves
vanishing and does not introduce a Gram coefficient inverse.
Monic contact division of these two functions and the conductors
above give a bounded Gram on $(R,S,U,T)$.
Also $G_0=xg_0(v)$ with bounded $g_0$, and exact contact gives a
double factor $(v-k)^2$, producing its bounded $U^2$ coefficient.

The full normals have zero-residue lifts on these roots.  For $L_X$,
write $L_X=uA_1+WB_1=-\kappa B_1S+u(A_1+\kappa B_1)$.
Both coefficient functions have identically zero leading part for
every $v$; exact normal contact makes $A_1+\kappa B_1$ vanish at
contact.  Monic division gives coefficients with factors $u\rho$,
$u^2$, or $uW$, and a $(v-k)$ quotient of an identically zero
leading function.  All have zero residue.  For $L_Y$, write
$L_Y=RB_2+\rho E_2$.  Its bounded even $E_2$ vanishes exactly at
contact, and both $B_2,E_2$ have identically zero leading part.
Monic division gives factors $x,\rho u,\rho W$, or the same
zero-leading-function quotient.  Thus the entire normal restoration
is bounded and has zero residue.

The soft residue has the form
\[
 q(k_0)R^2+\kappa_0^2c(k_0)S^2+A_1'RU+B_1'ST+E_1'U^2.
\]
It satisfies
\[
 E_1'>\frac{(A_1')^2}{4q(k_0)}+
          \frac{(B_1')^2}{4\kappa_0^2c(k_0)}.
\]
To prove this strict inequality, the leading scalar at
$u=\rho s,W=\rho w$ is
$\Phi=q s^2w^2+b s^2+c w^2-dsw+e$.
Complete its $w$ square:
\[
 \Phi=(qs^2+c)\left(w-\frac{ds}{2(qs^2+c)}\right)^2
           +\frac{\mathcal D(s^2,v^2)}{qs^2+c}.
\]
The global Laurent calculation gives a positive definite Hessian of
the second term in $(s,v)$ at $(1,k_0)$; both coordinate derivatives
are nonzero there.  The first square supplies the missing $w$
direction.  Hence the full three-variable Hessian is positive
definite.  On first variations $T=U$; the displayed inequality is
exactly its Schur complement after the invertible $(R,S)$ change.
Choose a real $\delta$ strictly between the second summand and
$E_1'-(A_1')^2/(4q(k_0))$.  Add $\delta$ times the exact zero
\[
 T^2-U^2-ST(v-k)-\kappa^{-1}RU(v-k)=0.
\]
The two blocks $(R,U)$ and $(S,T)$ are now strictly positive, and so
is the full six-root residue Gram.

If the ratio $W/Z$ is not infinitesimal, instead normalize $W=1$;
then $Z$ is bounded.  The earlier scalar is $aP^2$ plus a term of
positive residue $c(k_0)$; its normal lifts have zero residue on
$(P,1)$.  This is strict.  The case $Z=0$ is a literal sum of the
event squares and the three positive core squares.  The simultaneous
sign change $(u,v,X,Y)\mapsto(-u,-v,-X,-Y)$ covers $v_0=-k_0$.

\paragraph{The double-zero corner when $n_0=0$.}
Use instead the pure coefficient quantities
\[
 H_c=4n+4m+2,\quad \xi_c=4m/H_c,\quad
 \zeta=(4m+2)/H_c,\quad \eta_c=C/(C+A),
\]
and graph $X_g=\xi_cvW-\zeta uZ$, $Y_g=\eta_cuW$.
Put $\epsilon=4\alpha+4n(4m+2)/H_c>0$.
Then $\epsilon$ is infinitesimal, while $n/\epsilon$ and
$x/\epsilon$ are bounded.  Use the hard roots as before and scalar
roots
\[
 P=uW/\rho,\quad V=vW,\quad E=\rho Z,\quad
 B_\epsilon=\sqrt\epsilon\,uZ,\quad T=uvZ.
\]
Direct packet expansion gives exactly
\[
 f_{zz}=4\alpha+\epsilon u^2+4n\zeta^2u^4+4Dv^2
                         +u^2v^2 B_b,
\]
where $B_b$ is bounded with residue $2$.  Its first four terms
give coefficients $4\alpha/x$, $1+(4n\zeta^2/\epsilon)u^2$,
and $(4D/x)v^2$ on $E^2,B_\epsilon^2,E^2$ respectively.  The
$W^2$ contribution has bounded coefficients on $P^2,V^2$ with
positive residues $Q$ and $4m_0/(2m_0+1)$.  The $WZ$ contribution
is $uvWZ$ times a bounded function; use $uvWZ=vPE$ to give it zero
cross residue.  Thus the five scalar roots have strict diagonal
residue $(Q,4m_0/(2m_0+1),4E_*,1,2)$.

For completeness, $L_X$ has a zero-residue bounded coefficient on
$V$.  At $v=0$ its $Z$ coefficient is exactly $-8n\zeta u^3$;
on $B_\epsilon$ this is $-8n\zeta u^2/\sqrt\epsilon$, which is
infinitesimal because $n/\epsilon$ is bounded.  The remaining
$Z$ coefficient has a factor $uv^2$, giving a zero-residue multiple
of $T$.  For $L_Y$, the large-packet coefficient is now exactly
$B_2=8D_v(xA)\eta_cu^2$; the remaining terms lift as
\[
 L_Y=P(B_2+xe_w)+T(\rho u e_z)+E\,8v(D-\alpha),
\]
with bounded $e_w,e_z$.  All coefficients have zero residue.  The
hard residue is $\operatorname{diag}(4m_0+2,4(C_*+A_*))$.
This includes literal $n=0$ without deleting a coordinate.

\paragraph{The remaining commuting charts.}
The following exact core identities retain all unused event packets.
Every retained coefficient is nonnegative; in particular the $n$
packet can be retained even when $n=0$.  They prove strict bounded
Grams without dividing by a testing infinitesimal.

If $u$ is a finite unit and $v\to0$, take
\[
 R=(uW-Y)/\rho,\quad L=Y/\rho,\quad B=X+uZ,\quad V=vZ.
\]
Then $M=-B+(\rho v/u)(R+D_uL)$, and exactly
\begin{align*}
 H/2={}&(1+D_uv^2)B^2+u^2D_uV^2-2uvD_uBV\\
 &+2\rho R(uvB-D_uV)+xR^2-2D_u\rho LV+D_uxv^2L^2.
\end{align*}
Retain $n,D,\alpha$.  The residue on $(B,V,R,L)$ is diagonal
with entries $(4m_0+2,2u_0^2D_{u,0},4C_*,4A_*D_{u,0})$.

If $u\to\infty$ and $v$ is finite, put $p=1/u$ and use $p^2F$.
Take $R=(W-pY)/\rho$, $L=Y/\rho$, $B=Z+pX$, $V=vX$.
The homogenized $m$ amplitude is
$-B+\rho pvR+\rho vD_pL$, and the exact core is
\begin{align*}
 \widehat H/2={}&B^2+D_pV^2+2\rho D_pRV-2\rho pvRB+xR^2\\
 &+2D_p\rho pLV-2D_p\rho vLB+D_pxv^2L^2.
\end{align*}
Retain $n,D,\alpha$.  The residue on $(B,V,R,L)$ is diagonal
with entries $(4m_0+2,2,4C_*D_{v,0},4A_*D_{v,0})$.

If $v\to\infty$ and $u$ is bounded, put $q=1/v$ and use $q^2F$.
Take $R=(uW-Y)/\rho$, $L=Y/\rho$, $T=uqZ$, $N_1=\rho Z$.
Exactly,
\begin{align*}
 \widehat H/2={}&(D_u+q^2)X^2+D_uxL^2+T^2+2qXT
       +q^2xR^2+2uq\rho XR-2qRN_1-2D_uqLN_1,\\
 4D\widehat J^2={}&4(D/x)(-N_1+\rho uX+xqR)^2.
\end{align*}
Retain $n,m,\alpha$.  The residue on $(X,R,L,T,N_1)$ is diagonal
with entries
\[ (2D_{u,0},4C_*,4A_*D_{u,0},2,4D_*). \]

Finally if both angles are infinite, use $p=1/u,q=1/v$ and
$p^2q^2F$.  Set $R=(W-pY)/\rho$, $L=Y/\rho$, $T=qZ$.  Then
\[
 \widehat H/2=(1+p^2D_q)X^2+xq^2R^2+xD_pL^2+T^2
       +2\rho qXR+2pqXT-2\rho pRT-2\rho D_pLT.
\]
Retain $n,m,D,\alpha$.  The residue on $(X,R,L,T)$ is
diagonal with entries $(2,4C_*,4A_*,2)$.

\paragraph{Exhaustion and independence.}
These charts cover every commuting angle residue: the whole finite
$u=0$ axis, the finite-unit $v=0$ axis, either infinite axis, and
their intersections.  When both angle residues are noncommuting,
apply the inherited physical noncommuting-residue theorem at arbitrary
coefficient scales.  Complex states require no additional chart:
if $Z\ne0$, multiply by a phase and normalize $Z=1$; the real parent
is the sum of its real-state value and its imaginary-state value,
the latter having $Z=0$ and hence being a literal square sum.
If $Z=0$ this is already true of both parts.  Homogeneity covers all
projective state charts.  Roots may be unbounded throughout.
All inverses or radicals used in coefficients are pure coefficient
or pure testing operations on units, except $\sqrt\epsilon$ in
roots and its controlled coefficient ratio $n/\sqrt\epsilon$.
No mixed inverse and no comparison of independent infinitesimals
occurs.  Applying the strict bounded-tensor Gram lemma completes
the proof.
